\documentclass[letterpaper]{article}
\usepackage{amsmath,amssymb}
\usepackage[letterpaper,margin=1in]{geometry}
\setlength{\parindent}{0pt}
\newcommand{\noteheader}[1]{\par\bigskip\noindent\textbf{#1}\par\medskip}

\begin{document}

\textbf{CS 2050 --- Notes 09/21}

\noteheader{Resolution: simplifying a conclusion}

Suppose $p\lor q\lor r$ is true and each of $p$, $q$, and $r$ implies $s$.
Rewrite the implications as disjunctions, then resolve:
\[
\begin{array}{rll}
1.&p\lor q\lor r&\text{premise}\;\\[2pt]
2.&\neg p\lor s&\text{from }p\to s\;\\[2pt]
3.&\neg q\lor s&\text{from }q\to s\;\\[2pt]
4.&\neg r\lor s&\text{from }r\to s\;\\[4pt]
5.&q\lor r\lor s&\text{resolution of 1 and 2}\;\\[2pt]
6.&r\lor s&\text{resolution of 5 and 3}\;\\[2pt]
7.&s\lor s&\text{resolution of 6 and 4}\;\\[2pt]
8.&s&\text{idempotent law.}
\end{array}
\]

\noteheader{Another resolution example}

Simplify $(p\lor q)\land(\neg p\lor q)$:
\[
\begin{array}{rll}
1.&p\lor q&\text{from the conjunction}\;\\[2pt]
2.&\neg p\lor q&\text{from the conjunction}\;\\[2pt]
3.&q\lor q&\text{resolution of 1 and 2}\;\\[2pt]
4.&q&\text{idempotent law.}
\end{array}
\]

\clearpage

\noteheader{Exhaustive proof}

An exhaustive proof checks every possible case of a proposition. If the
proposition holds in every case, it is proved.

\noteheader{Proof by cases}

Use proof by cases when one argument does not cover every possibility.
Divide the possibilities into cases and prove the claim in each case.

\noteheader{Without loss of generality}

When cases are symmetric, it can be enough to prove one representative case
and state that the others follow by the same argument. This is called
proving the claim \textbf{without loss of generality} (WLOG).

\end{document}
